\documentclass{article}
\usepackage{amsmath, amssymb, amsthm}

\setlength{\parindent}{0pt}

\newtheorem{claim}{Claim}

\begin{document}

\begin{claim}
    Every planar graph has a node of degree at most 5.
\end{claim}

\begin{proof}
    Let G be a planar graph. Suppose for contradiction that all vertices of G 
    have degree at least 6. By the Handshaking Lemma:
    \[
    \sum_{v \in V} \deg(v) = 2|E|
    \]
    By our assumption, since every vertex has at least degree 6:
    \[
    2|E| = \sum_{v \in V} \deg(v) \geq 6|V|
    \]
    Simplifying, we get $2|E| \geq 6|V| \implies |E| \geq 3|V|$. Since G is planar, 
    we know that $|E| \leq 3|V| - 6$ from Euler's formula. However, $|E| \geq 3|V| 
    > 3|V| - 6 \geq |E| \implies |E| > |E|$, a contradiction.
    
    \medskip

    Hence every planar graph has a node of degree at most 5.
\end{proof}

\begin{claim}
    Any planar graph can be coloured with six colours.
\end{claim}

\begin{proof}
    We use induction on the number of vertices in the planar graph, which we denote 
    by $n$. Let $P(n)$ be the proposition that an $n$-vertex planar graph is 
    $6$-colourable.

    \medskip
    \textbf{Base case:}
    A $1$-vertex planar graph has degree $0$, hence it is $1$-colourable. Thus $P(1)$ 
    is true.

    \medskip
    \textbf{Inductive step:}
    Assume that $P(n)$ is true, and let $G$ be an $(n+1)$-vertex planar graph. By 
    Claim 1, $G$ has a node $v$ of degree at most $5$. Remove $v$ from $G$ and its
    incident edges, leaving an $n$-vertex planar graph $G'$. By the inductive 
    hypothesis, $G'$ is $6$-colourable. Now put $v$ back along with its incident 
    edges. Because $v$ has degree at most $5$, $v$ has at most $5$ adjacent vertices, 
    using at most $5$ colours, leaving at least $1$ colour for $v$. Hence $G$ is 
    $6$-colourable.

    \medskip

    Hence any planar graph can be coloured with six colours.
\end{proof}

\end{document}
